\documentclass{article}
\usepackage[preprint]{neurips_2026}

\usepackage[utf8]{inputenc}
\usepackage[T1]{fontenc}
\usepackage{hyperref}
\usepackage{url}
\usepackage{booktabs}
\usepackage{amsfonts}
\usepackage{amsmath}
\usepackage{microtype}
\usepackage{xcolor}

% Clickable links without the default boxes around citations/equation references.
\hypersetup{colorlinks=true, citecolor=black, linkcolor=black, urlcolor=blue}

% This is an ES 678 course paper, not a conference submission.
\makeatletter
\renewcommand{\@noticestring}{ES 678 Course Project}
\makeatother

\title{How Small Can a Reliable Prediction Set Be? Conformity-Score Design and
Class-Conditional Coverage in Conformal Classification}

\author{%
  Sumit Kumar Jha \\
  Roll Number: 26210180 \\
  ES 678 \\
  Indian Institute of Technology Gandhinagar \\
  \href{mailto:sumitjha@iitgn.ac.in}{\texttt{sumitjha@iitgn.ac.in}}
}

\begin{document}
\maketitle

\begin{abstract}
Conformal prediction wraps a fixed classifier so that it returns a \emph{set} of labels containing
the true class with probability at least $1-\alpha$, without retraining and without distributional
assumptions. This guarantee is only \emph{marginal}, however: it averages over the data and does not
control coverage for a particular class or a particular input. This project studies how the choice of
conformity score---including a proposed temperature-scaled, distance-augmented score---affects the
\emph{size} of prediction sets at a fixed guaranteed coverage, and audits the class-conditional
coverage that marginal guarantees fail to control, on MNIST and CIFAR-10 with an out-of-distribution
failure audit. The ES 678 depth component derives the coverage theorem and characterises the
efficiency--conditional-coverage frontier.
\end{abstract}

\section{Problem and Motivation}
\textbf{Question.} For a fixed, already-trained image classifier, how do (i) the choice of
\emph{conformity score} and (ii) the addition of feature-space \emph{distance} information affect
\emph{both} the \emph{size} of conformal prediction sets at a fixed, guaranteed coverage level
\emph{and} the \emph{class-conditional} coverage that marginal guarantees do not control?

\textbf{Intended prediction.} For each test input $x$ we output a \emph{set} of labels
$C(x)\subseteq\{1,\dots,K\}$ rather than a single label, such that $P\big(Y\in C(X)\big)\ge 1-\alpha$
without retraining and with no distributional assumptions beyond the exchangeability of calibration
and test data \citep{vovk2005algorithmic,shafer2008tutorial}. A prediction is made at inference time,
when no label is available and the input may be unfamiliar.

\textbf{What one data row represents.} One row is one image together with its ground-truth class
label; features are raw pixels and the model produces $\hat p(y\mid x)$. We evaluate set efficiency
and coverage, not raw accuracy, because a trustworthy system must know \emph{when} it is guessing.

\section{Data and Preparation}
All datasets are public and load directly through \texttt{torchvision}; there are no access or
sharing restrictions. One row is one image with its class label.

\begin{center}
\small
\renewcommand{\arraystretch}{1.2}
\begin{tabular}{p{2.8cm}p{3.9cm}p{5.5cm}}
\toprule
\textbf{Dataset} & \textbf{Role} & \textbf{Size / format} \\
\midrule
MNIST            & In-distribution (validity check) & 60{,}000 train / 10{,}000 test, $28\times28$ grayscale \\
CIFAR-10         & In-distribution (primary)        & 50{,}000 train / 10{,}000 test, $32\times32$ RGB \\
FashionMNIST, KMNIST & OOD for MNIST                & 10{,}000 test each \\
SVHN, CIFAR-100  & OOD for CIFAR-10                 & 10{,}000 or more test each \\
\bottomrule
\end{tabular}
\end{center}

\textbf{Access: yes} (fully public, downloaded by \texttt{torchvision}); the standard, unmodified
splits are used so results are comparable to the literature. \textbf{Quality checks.} We verify class
balance, pixel-value ranges, and near-duplicate images. \textbf{Preprocessing.} Inputs are normalised
with per-channel statistics computed on $D_{\text{train}}$ only; CIFAR-10 augmentation (random crop,
horizontal flip) is applied during training only, so no learned transformation touches the
calibration or test data. Because MNIST is nearly saturated (prediction sets collapse to singletons
at $\alpha\ge0.1$), the \emph{efficiency} study leads with CIFAR-10 and uses MNIST as a validity
check.

\section{Methods}
\textbf{Conformal construction.} Let $f$ be the classifier producing $\hat p(y\mid x)$, and let
$s(x,y)$ be a \emph{conformity score} (smaller $=$ more conforming). Given the $n$ calibration scores
$S_1,\dots,S_n$, split conformal takes the quantile
\begin{equation}
  \hat q = S_{(k)}, \qquad k=\big\lceil (n+1)(1-\alpha)\big\rceil,
  \label{eq:quantile}
\end{equation}
where $S_{(1)}\le\dots\le S_{(n)}$ are the sorted calibration scores, and returns the set
\begin{equation}
  C(x)=\{\,y : s(x,y)\le \hat q\,\}.
  \label{eq:set}
\end{equation}
Under exchangeability this guarantees, for finite $n$,
\begin{equation}
  1-\alpha \;\le\; P\big(Y\in C(X)\big) \;\le\; 1-\alpha+\frac{1}{n+1}.
  \label{eq:bound}
\end{equation}

\textbf{Baselines.} (1) Top-1 point prediction (accuracy reference: no set and no guarantee);
(2) a simple logistic-regression classifier on MNIST, to justify the use of a CNN; and (3) the
\emph{naive} conformity score $s(x,y)=1-\hat p(y\mid x)$, which is valid but not adaptive.

\textbf{Candidate scores.} The first three are established; the last is ours.
\begin{align}
  s_{\text{naive}}(x,y)   &= 1-\hat p(y\mid x), \label{eq:naive}\\[2pt]
  s_{\text{APS}}(x,y)     &= \sum_{j<k_y}\hat p_{(j)}(x)+u\,\hat p_{(k_y)}(x),
                             \qquad u\sim\mathrm{Unif}(0,1), \label{eq:aps}\\[2pt]
  s_{\text{RAPS}}(x,y)    &= s_{\text{APS}}(x,y)+\lambda\,(k_y-k_{\text{reg}})^{+},
                             \label{eq:raps}\\[2pt]
  s_\gamma(x,y)           &= \big(1-\hat p_T(y\mid x)\big)+\gamma\,\tilde d_M(x,y),
                             \qquad \gamma\ge 0, \label{eq:proposed}
\end{align}
with the temperature-scaled probabilities and the Mahalanobis distance defined by
\begin{align}
  \hat p_T(y\mid x) &= \operatorname{softmax}\!\big(z(x)/T\big), \label{eq:temp}\\[2pt]
  d_M(x,y)          &= \sqrt{\big(\phi(x)-\mu_y\big)^{\!\top}\,\Sigma^{-1}\,\big(\phi(x)-\mu_y\big)},
                       \label{eq:maha}
\end{align}
where $z(x)$ are the logits, $\phi(x)$ the penultimate features, and $\mu_y,\Sigma$ the per-class
means and a shared covariance. Here $k_y$ is the rank of $y$ in the sorted probabilities, APS is
Adaptive Prediction Sets \citep{romano2020classification}, RAPS adds a rank penalty
\citep{angelopoulos2021learn}, $T$ is fit by minimising the negative log-likelihood
\citep{guo2017calibration}, and the distance term follows the Mahalanobis OOD score
\citep{lee2018simple}.

\textbf{Where every quantity is fit (leakage discipline).} The classifier is trained on
$D_{\text{train}}$. The temperature $T$, the RAPS parameters $\lambda,k_{\text{reg}}$, our distance
weight $\gamma$, and method selection all use $D_{\text{design}}$ only. The class Gaussians
$(\mu_y,\Sigma)$ and the standardisation of $d_M$ are also fit on $D_{\text{design}}$ (held out from
training), so they are not in-sample for the classifier. The quantile $\hat q$ of \eqref{eq:quantile}
uses $D_{\text{cal}}$ only, and $D_{\text{test}}$ is opened once, for the final table. Mondrian
(class-conditional) conformal instead calibrates a separate per-class quantile $\hat q_c$ from
calibration points with label $c$, guaranteeing $P(Y\in C(X)\mid Y=c)\ge 1-\alpha$
\citep{vovk2012conditional}. Our score \eqref{eq:proposed} reduces to the temperature-scaled baseline
at $\gamma=0$, giving a clean ablation.

\section{Evaluation Design}
\textbf{Split.} A fixed, disjoint four-way split is used \emph{before} any comparison:
$D_{\text{train}}$ (fit classifier) / $D_{\text{design}}$ (early stopping; fit
$T,\gamma,\lambda,k_{\text{reg}}$; method selection) / $D_{\text{cal}}$ (compute $\hat q$ only) /
$D_{\text{test}}$ (final metrics). Model selection and tuning are thus kept entirely separate from
final testing.

\textbf{Metrics.} We report (i) \emph{marginal empirical coverage}, read against the finite-sample
window in \eqref{eq:bound} so that a result counts as valid only if the observed coverage lies inside
it; (ii) \emph{mean}, median and 90th-percentile \emph{set size}---the efficiency metric, because
smaller sets are strictly more useful at the same error budget \citep{sadinle2019least}; (iii)
singleton and empty-set rates; (iv) \emph{class-conditional} coverage and the worst-class gap; (v)
subgroup coverage by confidence decile; and (vi) coverage and size on the OOD sets. The study is
repeated for $\alpha\in\{0.01,0.05,0.10,0.20\}$ and reported as mean $\pm$ standard deviation over
3--5 seeds.

\section{Results and Error Analysis}
Because this is a proposal, we state the expected results as testable hypotheses; a negative result
is acceptable and will be reported as such.
\begin{itemize}
  \item \textbf{H1 (validity).} Empirical marginal coverage lies in the window of \eqref{eq:bound}
        for every score. \emph{Falsifier:} coverage outside the window.
  \item \textbf{H2 (efficiency).} APS yields a smaller mean set size than the naive score at matched
        coverage. \emph{Falsifier:} naive $\le$ APS.
  \item \textbf{H3 (calibration).} Temperature scaling reduces the APS set size on the miscalibrated
        CIFAR-10 classifier. \emph{Falsifier:} no improvement.
  \item \textbf{H4 (distance term).} $s_\gamma$ with $\gamma>0$ reduces the worst-class coverage gap
        versus temperature-only, at a bounded cost in mean set size. \emph{Falsifier:} the gap is
        unchanged or worse.
  \item \textbf{H5 (conditional coverage).} Mondrian attains the smallest worst-class gap but the
        largest mean set size. \emph{Falsifier:} Mondrian is not best on conditional coverage.
\end{itemize}
\textbf{Error analysis.} We will inspect which classes and inputs break coverage, whether failures
concentrate in low-confidence deciles, and how coverage and size behave on the OOD sets where
exchangeability fails.

\section{Graduate Depth Analysis}
The primary reference discussed is the adaptive prediction sets of
\citet{romano2020classification}, extended with rank regularisation \citep{angelopoulos2021learn},
calibration
\citep{guo2017calibration}, and a Mahalanobis feature-space distance \citep{lee2018simple}. The PG
depth deliverable is a study of the \emph{efficiency--conditional-coverage frontier}, containing:
\begin{enumerate}
  \item \textbf{Derivation / theory.} A full proof, under exchangeability, of the split-conformal
        coverage bound \eqref{eq:bound}, with the finite-sample slack $1/(n+1)$ interpreted
        explicitly.
  \item \textbf{Calibration / uncertainty analysis.} Class-conditional coverage per class, the
        worst-class coverage gap, and a comparison of the proposed $s_\gamma$ against Mondrian
        (guaranteed) and temperature-only baselines.
  \item \textbf{Sensitivity / ablation.} An $\alpha\times\gamma$ grid; the ablation
        $s_{\text{APS}}\to s_{\text{APS}}(T)\to s_{\text{APS}}(T)+\text{distance}$ that isolates the
        contribution of each term.
  \item \textbf{Failure-mode audit.} Which classes and inputs break coverage; behaviour on OOD sets
        where exchangeability fails and the guarantee is void (see covariate-shift conformal
        \citep{tibshirani2019conformal}); and the effect of a small calibration set $n$ on the slack
        in \eqref{eq:bound}.
\end{enumerate}
This satisfies four of the addendum's listed options (derivation, calibration/uncertainty,
sensitivity/ablation, failure-mode audit).

\section{Limitations and Conclusion}
\textbf{One likely challenge.} \emph{Conditional coverage is not achievable for arbitrary inputs}
(a known impossibility), so validity must be framed as \emph{marginal} coverage plus an
\emph{empirical} conditional-coverage audit. A related risk is \emph{leakage}: tuning the conformal
score on the same split used to compute $\hat q$ breaks exchangeability and silently invalidates the
guarantee---hence the strict $D_{\text{design}}\,/\,D_{\text{cal}}$ separation above. A third risk is
that the proposed distance term yields only a small or null gain; this is anticipated (H4) and will
be reported honestly rather than overclaimed. \textbf{Mitigations:} separate design and calibration
splits; multiple seeds with error bars; a disclosed fallback CIFAR-10 backbone; and leading with the
CIFAR-10 efficiency study while treating MNIST as a validity check. \textbf{Conclusion.} By holding
the classifier fixed and varying only the conformity score, the project answers a precise question:
how to shrink prediction sets without losing the coverage guarantee, and where that guarantee fails
conditionally.

\section{Reproducibility and AI-Use Disclosure}
The code is organised as a pipeline (\texttt{data}$\to$\texttt{model}$\to$\texttt{calibrate}$\to$
\texttt{scores}$\to$\texttt{conformal}$\to$\texttt{evaluate}) with a single entry point that
regenerates every figure and table from YAML configs and fixed seeds; all datasets load via
\texttt{torchvision}. \textbf{Compute.} The classifier is trained once (three seeds for stability);
the reported error bars vary only the calibration/test split and the randomised-APS seed, so the full
grid is feasible on an Apple-silicon laptop (MPS). \textbf{Timeline (four weeks).} Week 1: data,
four-way split, and a naive-conformal coverage sanity check. Week 2: APS/RAPS and temperature
scaling. Week 3: the distance-augmented score plus the conditional-coverage and OOD audits. Week 4:
the four-page paper, the reproduction script, and defence preparation. \textbf{Tools and AI use.}
Python 3.11, PyTorch (MPS on Apple Silicon), torchvision, NumPy, SciPy, pandas, Matplotlib; \LaTeX{}
for this proposal. AI tools were used for brainstorming, boilerplate and editing; every equation,
number and claim is verified independently by the author and disclosed in the final paper.

\begin{thebibliography}{9}
\bibitem[Angelopoulos et al.(2021)]{angelopoulos2021learn}
Anastasios N. Angelopoulos, Stephen Bates, Michael I. Jordan, and Jitendra Malik.
\newblock Uncertainty sets for image classifiers using conformal prediction.
\newblock In \emph{International Conference on Learning Representations}, 2021.

\bibitem[Guo et al.(2017)]{guo2017calibration}
Chuan Guo, Geoff Pleiss, Yu Sun, and Kilian Q. Weinberger.
\newblock On calibration of modern neural networks.
\newblock In \emph{International Conference on Machine Learning}, 2017.

\bibitem[Lee et al.(2018)]{lee2018simple}
Kimin Lee, Kibok Lee, Honglak Lee, and Jinwoo Shin.
\newblock A simple unified framework for detecting out-of-distribution samples and adversarial
attacks.
\newblock In \emph{Advances in Neural Information Processing Systems}, 2018.

\bibitem[Romano et al.(2020)]{romano2020classification}
Yaniv Romano, Matteo Sesia, and Emmanuel J. Cand\`es.
\newblock Classification with valid and adaptive coverage.
\newblock In \emph{Advances in Neural Information Processing Systems}, 2020.

\bibitem[Sadinle et al.(2019)]{sadinle2019least}
Mauricio Sadinle, Jing Lei, and Larry Wasserman.
\newblock Least ambiguous set-valued classifiers with bounded error levels.
\newblock \emph{Journal of the American Statistical Association}, 114(525):223--234, 2019.

\bibitem[Shafer and Vovk(2008)]{shafer2008tutorial}
Glenn Shafer and Vladimir Vovk.
\newblock A tutorial on conformal prediction.
\newblock \emph{Journal of Machine Learning Research}, 9:371--421, 2008.

\bibitem[Tibshirani et al.(2019)]{tibshirani2019conformal}
Ryan J. Tibshirani, Rina Foygel Barber, Emmanuel J. Cand\`es, and Aaditya Ramdas.
\newblock Conformal prediction under covariate shift.
\newblock In \emph{Advances in Neural Information Processing Systems}, 2019.

\bibitem[Vovk(2012)]{vovk2012conditional}
Vladimir Vovk.
\newblock Conditional validity of inductive conformal predictors.
\newblock In \emph{Asian Conference on Machine Learning}, 2012.

\bibitem[Vovk et al.(2005)]{vovk2005algorithmic}
Vladimir Vovk, Alexander Gammerman, and Glenn Shafer.
\newblock \emph{Algorithmic Learning in a Random World}.
\newblock Springer, 2005.
\end{thebibliography}

\end{document}
